\section{Existence of the flat deformation and its mass comparison}
\label{sec:fc-existence}

We retain the prepared exterior $\Omega$, the inner faces
$B=B_b\sqcup B_w$, and the four-stage family of
Proposition~\ref{prop:fl-homotopies}.  The parameter of that family is
$\xi\in[0,4]$, with $\xi=0$ the physical system.  All the geometric
choices in Lemma~\ref{lem:fl-threshold-preparation} are fixed throughout
this section.  A finite truncation is denoted by $\Omega_R$.
The normal $\nu$ on $B$ points into $\Omega_R$; its negative is the
normal used in the divergence theorem.  We write
\[
 \sigma=|df|_g,\qquad e=\nabla f/\sigma\quad(\sigma>0),\qquad
 |\zeta|_{A_j}^2=A_j(\zeta,\zeta)
\]
for a covector $\zeta$.  Expressions at $\sigma=0$ are understood by
their smooth extensions established in Section~\ref{sec:fl-system-floor}.

We first obtain collar and global bounds for arbitrary smooth solutions.
We then construct a solution by scalar solvability, degree, and exhaustion,
and conclude with the curvature and mass comparison.

There are two distinct types of constants.  A symbol $P_N$ denotes a
quantity bounded by $C(1+N)^q$, where $C,q$ depend on the fixed data,
collars, and $\epsilon$, but not on $N$, $R$, or $\xi$.  Such a
polynomial may be increased at successive occurrences.  A constant
$C_N$ may depend arbitrarily on the fixed positive parameters
$\epsilon,N,\ell,\tau$, but remains independent of $R$ and $\xi$.
In particular,
\begin{equation}\label{eq:fc-parameter-ledger}
 \ell=e^{-\epsilon N},\qquad \tau=\ell^{3/2},\qquad
 \ell\le l\le e,\qquad
 \frac{4d}{4+N}\le\chi\le d,\qquad
 P_N\frac{\tau}{\ell}\longrightarrow0.
\end{equation}
All estimates preceding the existence proof concern arbitrary smooth
solutions with $t\ge-\epsilon$, including a putative solution touching
that level.  No estimate for a Hessian or upper bound for $Z$ is
assumed at this stage.

\subsection{Closed sublevels and signed collar estimates}

Write $\mathcal B_c$ for the full black trapped region at threshold
$c$, and $\mathcal W_c$ for the corresponding white region with
$\mathcal B_{c_b}$ held fixed.  These regions, including their
Hausdorff right continuity at the selected thresholds, are those of
Lemma~\ref{lem:fl-threshold-preparation}.

\begin{lemma}[Separation from the opposite height]
\label{lem:fc-compact-separation}
In the first two stages, if a fixed compact subset $Q$ of the closed
prepared exterior is disjoint from $\mathcal B_{c_b}$, there are
$\eta_Q>0$ and $N_Q$ such that
\[
 h\ge b_-+\eta_Q\quad\hbox{on }Q
\]
for $N\ge N_Q$ and all sufficiently large allowed truncations.
If $Q$ is disjoint from $\mathcal W_{c_w}$, then instead
$h\le b_+-\eta_Q$.  The compact sets may meet a face of the opposite
color.  These assertions are uniform over the two stage parameters.
\end{lemma}

\begin{proof}
Consider any sequence $N_i\to\infty$, $R_i\to\infty$ of the
indicated solutions.  Lemma~\ref{lem:fl-height} gives, with fixed
constants,
\begin{equation}\label{eq:fc-height-sequence}
 |h_i|\le C_0,
 \qquad |h_i|\le C(r^{-\gamma}-R_i^{-\gamma})
 \quad\hbox{on the end},\qquad \frac34<\gamma<1.
\end{equation}
Extend each $h_i$ across every white face by the constant $b_+$,
and extend the fixed metric and tensor smoothly into those short
collars.  These extensions of $h_i$ are continuous; no uniform
boundary modulus is required.  Let
\[
 \underline h(x)=
 \lim_{j\to\infty}\inf\{h_i(y):i\ge j,
                         \operatorname{dist}(x,y)<j^{-1}\}.
\]
This is a bounded lower semicontinuous function satisfying the
limiting end bound.

Choose $b_-<k'<0$ so near $b_-$ that the stage-two collar term is
nonpositive at every height at most $k'$.  Its positive white part
vanishes in this fixed low-height range.  Suppose $\phi$ is a smooth
lower test of $\underline h$ at $x$, where
$\underline h(x)\le k'$ and $d\phi(x)\ne0$.  For now exclude the
black faces.  Subtracting a small positive multiple of
$\operatorname{dist}(x,\cdot)^4$ makes contact strict without
altering the two-jet.  Minimization on a small closed ball supplies
points $x_i\to x$ and constants tending to zero for which the
translated tests touch $h_i$ from below and
$h_i(x_i)\to\underline h(x)$.  They cannot touch the boundary of
the small ball.  They cannot touch a white face or its extended
side either, because $h_i=b_+>k'$ there.  Thus these are interior
contacts, even if $x$ itself lies on a white face.

At these contacts $df_i=\tau_i^{-1}d\phi$, whence
\[
 d_i\le\frac{\tau_i^2}{\ell_i^2|d\phi(x_i)|^2}\longrightarrow0,
 \qquad a_i\longrightarrow1,\qquad\chi_i\longrightarrow0.
\]
The positive Hessian matrix in the trace equation gives
\[
 \tr_{A_{\chi_i}}K+
 \frac{l_i}{\tau_i\sqrt{D_i}}
       \tr_{A_{\chi_i}}\Hess_g\phi
 \le F_i(x_i).
\]
Here $l_i/(\tau_i\sqrt{D_i})=a_i/|d\phi|$.  The bound on the
collar addition and $C\le C_b$ therefore give
\begin{equation}\label{eq:fc-test-expansion}
 \frac{\tr_{(d\phi)^\perp}\Hess_g\phi}{|d\phi|}
       +\tr_{(d\phi)^\perp}K\le k'+C_b.
\end{equation}
The left side is the expansion of the level set, oriented toward
increasing $\phi$.  Removing the fourth-order localization proves
the same assertion for a non-strict lower test.

To pass from functions to all supports of a closed set, fix
$b_-<k<k'$ and put $E_k=\{\underline h\le k\}$.  Define
\[
 q_j(s)=\bigl(1+\exp[-j^2(s-k-j^{-1})]\bigr)^{-1},
 \qquad v_j=q_j(\underline h).
\]
The lower relaxed limit of $v_j$ is zero on $E_k$ and one off it.
At an equality point this follows from
$q_j(k)=(1+e^j)^{-1}\to0$; off $E_k$, lower semicontinuity gives
a neighborhood whose heights exceed $k$ by a fixed positive amount.
Let $\psi$ be a nonzero-gradient lower test of this limit at a zero
value.  Strict localization gives lower tests of $v_j$ at points
$y_j\to x$ with $v_j(y_j)\to0$.  Their original heights are below
$k'$ for large $j$, since $q_j(k')\to1$.  For each fixed such $j$
the smooth inverse $q_j^{-1}$ turns this into a lower test of
$\underline h$ with finite nonzero gradient.  Apply
\eqref{eq:fc-test-expansion} to that inverse test, taking the
solution-sequence limit with $j$ fixed, and only then let
$j\to\infty$.  An increasing change of defining function preserves
the level expansion: the additional Hessian term is a multiple of
$d\psi\otimes d\psi$ and has zero tangential trace.  Consequently
every such test $\psi$ satisfies
\eqref{eq:fc-test-expansion}.

A smooth local exterior support of $E_k$ has a defining function
that is zero at contact and nonpositive on $E_k$.  Shrinking and
scaling the neighborhood so its positive values are below one
makes it exactly a lower test of the preceding zero--one function.
Thus every exterior support has expansion at most $k'+C_b$.
In particular $E_k$ cannot touch a white face: the reversed face is
an exterior support with expansion
\[
 -H+\tr_TK=-c_w>0,
\]
whereas $k'+C_b<0$ after $k'$ is chosen sufficiently close to $b_-$.
Equivalently, in a signed white collar with $s\ge0$ in $\Omega$,
a small multiple of $-s$ is the offending lower test.  This rules
out boundary layers without assuming that the Dirichlet values
persist in the relaxed limit.  The end bound makes $E_k$ compact.

Adjoin the entire black region and fill complementary components
which contain no designated outer barrier.  An exterior support at
a black frontier point also supports the smooth black region, so
its expansion is at most $c_b<k'+C_b$.  At other frontier points
the preceding argument applies.  Filling produces no frontier
point outside the old closed set; every support of the larger
filled set is still a support of that old set.  The support bound
is therefore preserved.  The original required inner collars are
covered by the black region, and the outer barriers remain outside.
Lemma~\ref{lem:fl-support-enclosure} places the resulting set in
$\mathcal B_{c''}$ for any sufficiently close threshold
$c''>k'+C_b$.

For a fixed compact $Q$ disjoint from $\mathcal B_{c_b}$, right
continuity gives $\delta>0$ with
$\mathcal B_{c_b+\delta}\cap Q=\varnothing$.  Choose
\[
 b_-<k<k',\qquad k'+C_b<c''<c_b+\delta.
\]
Then $E_k\cap Q=\varnothing$ for every sequence considered above.
Failure of a uniform lower margin would give violating points in
$Q$, a convergent subsequence, and a relaxed limit value at most
$k$ there.  This contradiction proves the lower assertion.
Apply the argument to $(-h,-K,-C)$ for the upper assertion.  In
the white problem, filling retains all reversed black outer faces.
If the white region is empty, a nonempty violating sublevel would
still provide the seed for Lemma~\ref{lem:fl-support-enclosure};
right continuity at the empty region rules out such seeds at
arbitrarily close higher thresholds.  Multiple components and
bounded complementary pockets cause no change.

For the truncation quantifier start with any $R_{\min}(N)\to\infty$
satisfying the end estimates, and enlarge it as needed.  Failure
for arbitrarily large $N$ and $R\ge R_{\min}(N)$ would provide
exactly a sequence excluded above.  Only finitely many fixed
compact sets will be used below, so their thresholds can be
chosen simultaneously.
\end{proof}

\begin{lemma}[Collar slopes]\label{lem:fc-collar-slopes}
For sufficiently large $N$, the first two stages satisfy, with
fixed $c,C_1>0$,
\begin{equation}\label{eq:fc-signed-slopes}
 \frac c\tau\le\partial_\nu f\le\frac{C_1}\tau\quad\hbox{on }B_b,
 \qquad
 \frac c\tau\le-\partial_\nu f\le\frac{C_1}\tau\quad\hbox{on }B_w.
\end{equation}
In the third stage $|\partial_\nu f|\le C_1/\tau$.  In the first
two stages $b_-\le h\le b_+$ on every fixed compact region needed
for the scalar-curvature comparison.
\end{lemma}

\begin{proof}
Let $s$ be a fixed outward leaf parameter in a black collar, and
let $n_s=\nabla s/|ds|$.  For $h_0=b_-+\psi(s)$ with $\psi'>0$,
evaluate the operator at a comparison contact, where the gradients
of the test and solution agree.  The floor yields
\begin{equation}\label{eq:fc-collar-smallness}
 d\le C\frac{\tau^2}{\ell^2(\psi')^2},\qquad
 1-a\le d,\qquad N\chi\ge c_1d\quad(N\ge4).
\end{equation}
The trace on the test is exactly
\begin{equation}\label{eq:fc-collar-trace}
 P_s+aH_s+\chi K(n_s,n_s)
 +a\chi\left(
       \frac{\Hess s(n_s,n_s)}{|ds|}
                   +|ds|\frac{\psi''}{\psi'}\right),
\end{equation}
where $P_s=\tr_{T_s}K$.  In particular it is the leaf expansion
$H_s+P_s$, an $O(d)$ error, and the displayed logarithmic-slope
term.  This calculation differentiates neither the input $Z$ nor
an unknown coefficient.  For a negative slope the Hessian terms
have $a$ replaced by $-a$.

Since $H_s+P_s\ge c_b$ and $C=C_b-k_Bs$, choose small positive
slopes with $\psi(0)=0$ and $\psi(s)\le k_Bs/2$.
Lemma~\ref{lem:fc-compact-separation} orders the test below the
solution at the outer end of a fixed short collar.  The collar
addition vanishes in this compatible direction.  The residual of
the trace equation is at least
\[
 \frac{k_Bs}{2}-Cd+a\chi|ds|\frac{\psi''}{\psi'}.
\]
Choose a nonnegative smooth $\beta_N$ equal to $AN$ on
$[0,N^{-1}]$, supported in $[0,2N^{-1}]$, and with integral at most
$2A$.  Set
\[
 \psi'(s)=m\exp\!\left(\int_0^s\beta_N(q)\,dq\right).
\]
A fixed large $A$ makes the final term dominate $Cd$ on the first
interval, by \eqref{eq:fc-collar-smallness}.  Beyond that interval
$k_Bs/2$ dominates $Cd\le C\ell$ for large $N$.  A sufficiently
small fixed $m>0$ ensures all height and outer-edge requirements,
since the slope multiplier is at most $e^{2A}$.  This is a strict
lower barrier with its $h$-slope in a fixed positive interval.

For an upper barrier use
$\psi'(s)=M\exp(-\int_0^s\beta_N)$.
A fixed large $M$ makes its height gain dominate the outer-edge
height bound and all $O(s)$ variations of the leaf expansion and
offset.  Its logarithmic-slope term has the negative sign and
dominates the $O(d)$ errors initially; the linear height margin
dominates them thereafter.  At a negative minimum of $h-h_0$ the
Hessian of $h$ dominates that of $h_0$, their gradients and $Z$
agree, and their implicit $t,l,A_\chi$ agree.  Strict height
monotonicity contradicts a positive lower-barrier residual.
At a positive maximum the signs reverse.  This proves comparison
and the black estimates.  The sign-reversed construction proves
the white estimates.

In stage three start at its prescribed face value and use a large
positive-slope upper barrier and a large negative-slope lower
barrier, both with logarithmic absolute-slope derivative
$-\beta_N$.  The directional inequalities in
Proposition~\ref{prop:fl-homotopies} give the initial residual
signs.  Formula~\eqref{eq:fl-homotopy-three} evaluates its shifted
height at the original face value, so these signs persist along
the interpolation.  At the unshifted face height all errors are
$O(s+d)$; changing the height supplies a margin at least the
absolute height change.  Large fixed slopes absorb the $O(s)$
errors and order the outer edge, and the same logarithmic
correction handles the initial $O(d)$ errors.  This proves the
absolute upper slope bound without using separation in stage
three.  Finally combine the own-color lower collar bounds with
Lemma~\ref{lem:fc-compact-separation} on the rest of a fixed
compact set.  The opposite bound on each collar follows from that
lemma as well.  A finite cover proves the asserted height range.
\end{proof}

\begin{lemma}[Collar trace inequalities]
\label{lem:fc-trace-estimate}
Let $\mathcal C$ be a fixed union of shortened inner collars.
In the first two stages, for every nonnegative smooth $\phi$,
\begin{align}
 \int_B L\phi\,dA_g
 &\le P_N\int_{\mathcal C}u
       \bigl((1+\sqrt{\mathcal T})\phi+|d\phi|_{A_\chi}\bigr)
            \,dV_g,\label{eq:fc-weighted-trace}\\
 \int_B \phi\,dA_g
 &\le P_N\int_{\mathcal C}\sqrt D
       \bigl((1+\sqrt{\mathcal T})\phi+|d\phi|_{A_\chi}\bigr)
            \,dV_g.\label{eq:fc-unweighted-trace}
\end{align}
Fixed collar cutoffs equal to one on $B$ are understood; their
derivatives are included in the coefficient of $\phi$.
\end{lemma}

\begin{proof}
Set $W=l\bar\nabla f=l\nabla f/D=\sqrt d\,a e$.
Its $\bar g$-norm is at most one, and $uW=Lw$.  For a covector
$\zeta$,
\[
 |\zeta(W)|\le|\zeta|_{A_d}
       \le(1+N/4)^{1/2}|\zeta|_{A_\chi}.
\]
Direct differentiation gives
\begin{align*}
 u^{-1}\operatorname{div}_g(uW)
 &=\sqrt d\bigl(\tr_{A_d}H^f+(dp+p+2)a\,\partial_e t\bigr),\\
 D^{-1/2}\operatorname{div}_g(\sqrt D\,W)
 &=\sqrt d\bigl(\tr_{A_d}H^f+dp\,a\,\partial_e t\bigr).
\end{align*}
Both are bounded by $P_N(1+\sqrt{\mathcal T})$ using
Lemma~\ref{lem:fl-alg-coercivity}.  In particular their axial
Hessian coefficient is $d^{3/2}$, and their axial $dt$ coefficient
is at most $(2N+2)\sqrt d$; the other Hessian terms are transverse.
No factor $\ell^{-1}$ occurs.  At a black or white face, respectively,
$uW_\nu=La$ or $-La$, and $a\ge1/2$ by
\eqref{eq:fc-signed-slopes}.  Apply the divergence theorem with
the corresponding sign, the collar cutoff, and $\phi$.  Replacing
$u$ by $\sqrt D$ gives the second inequality.  The formulas also
hold at zero gradients by continuity.
\end{proof}

\subsection{Global bounds before regularity}

\begin{lemma}[High-level weighted energy]\label{lem:fc-high-energy}
In stage one there is $M_0=M_0(N)>0$ such that, if $j_0$ is a
smooth increasing cutoff equal to zero below $M_0$ and to one
above $M_0+1$, with bounded derivative, then
\begin{equation}\label{eq:fc-high-energy}
 \int_{\Omega_R}u j_0
       (\mathcal T+\rho+\tau a\sigma)\,dV_g
 +\int_{\Omega_R}u j_0'|dZ|_{A_\chi}^2\,dV_g
 \le C_N.
\end{equation}
Consequently, on the ordinary product graph $\Gamma$ of $f$,
every fixed compact base patch $Q$ satisfies
\begin{equation}\label{eq:fc-graph-L2}
 \int_{\Gamma|Q}e^{2Z}\,dV_\Gamma\le C_{N,Q}.
\end{equation}
\end{lemma}

\begin{proof}
On the penalty band $-\epsilon\le t\le-\epsilon+N^{-1}$,
$l\asymp\ell$ and $D=e^{8(Z-t)}$.  Thus $a\sigma$ tends uniformly
to infinity as $Z\to\infty$ on that band.  The penalty weights
are bounded on the entire exterior, so a sufficiently large
$M_0(N)$ ensures
\begin{equation}\label{eq:fc-high-source}
 \Xi\ge\delta_0\mathcal T+\rho
                    +\tfrac12\delta_0\tau a\sigma
                 \quad\hbox{on }\{Z>M_0\}.
\end{equation}
Outside the penalty band this is immediate.  Testing the first
stage divergence equation by $j_0(Z)$ has no outer boundary term,
since the outer value of $Z$ is zero.  The drift term is bounded by
\[
 u j_0'|\langle dZ,K(w,\cdot)\rangle_{A_\chi}|
 \le 2u j_0'|dZ|_{A_\chi}^2+C u j_0'|K|^2.
\]
The latter integral is bounded by $C_N$: the floor and the bounded
$Z$ strip bound $u$ there, while $K=O(r^{-2})$ has finite squared
integral on the three-dimensional AF end.

At either color of face, the matching signs give
$-H+w_\nu(F-P_B)=-(1-a)H$.
The omitted-drift term has normal component bounded by
$C\chi\le Cd$, and the scalar boundary multiplier lies in
$[0,1]$.  Hence, uniformly in the first stage,
\begin{equation}\label{eq:fc-boundary-flux-small}
 |(V_\xi)_\nu|\le Cud=CL\sqrt d
                  \le C\frac\tau\ell L.
\end{equation}
Integration and \eqref{eq:fc-high-source} give the left side of
\eqref{eq:fc-high-energy}, up to fixed positive coefficients,
bounded by $C_N+C(\tau/\ell)\int_B L j_0$.
Use \eqref{eq:fc-weighted-trace}.  On the fixed collars the
positive minimum of $\rho$ implies
\[
 1+\sqrt{\mathcal T}\le C(\rho+\delta_0\mathcal T).
\]
Thus the terms with $j_0$ absorb for $N$ large, by
\eqref{eq:fc-parameter-ledger}.  The derivative term is bounded by
\[
 u j_0'|dZ|_{A_\chi}^2
                  +C(P_N\tau/\ell)^2u j_0'.
\]
Its second summand has a bounded integral on the compact collar
and the bounded $Z$ strip.  This proves
\eqref{eq:fc-high-energy}.

Now constants may have arbitrary fixed-$N$ dependence.  The
measures $dV_\Gamma=\sqrt{1+\sigma^2}\,dV_g$ and
$\sqrt D\,dV_g$ are comparable because $\ell\le l\le e$.
Moreover
\[
 e^{2Z}=e^{2t}D^{1/4},\qquad
 D^{1/4}\le C_N(1+\tau a\sigma),\qquad u=le^{2t}\sqrt D.
\]
The middle bound follows by separating $l\sigma\le1$ and
$l\sigma>1$, on which $a\ge1/\sqrt2$.
On a fixed compact $Q$, its positive minimum of $\rho$ and
\eqref{eq:fc-high-energy} control the high-level part of
\eqref{eq:fc-graph-L2}.  On $Z\le M_0+1$ the floor bounds $D$,
$\sigma$, and $e^{2Z}$, so the remaining graph volume is bounded
by a constant times the fixed base volume.  This is a local
weighted estimate, with no global graph-volume assumption.
\end{proof}

\begin{lemma}[Graph Sobolev Inequality and iteration]
\label{lem:fc-graph-iteration}
In stage one, $Z$ has a uniform upper bound on every fixed compact
subset of the closed prepared exterior.  The bound is independent
of the truncation and first-stage parameter at fixed $N$.
\end{lemma}

\begin{proof}
Put $d_o=(1+\sigma^2)^{-1}$.  The matrices $A_{d_o},A_d,A_\chi$
are uniformly comparable at fixed $N$, without an upper bound on
$Z$.  Here $|\nabla_\Gamma\omega|=|d\omega|_{A_{d_o}}$.
Extend a fixed base patch to a compact smooth enlargement and
isometrically embed it in Euclidean space by Nash's Theorem
\cite{nash1956}.  Its product with the line immerses the ordinary
graph isometrically.  The base embedding contributes a bounded
amount to the graph mean curvature vector.  Its mean curvature
in the product is
\[
 H_\Gamma=
 \frac{\sqrt D}{l\sqrt{1+\sigma^2}}
       \bigl(\tr_{e^\perp}H^f+d_o H^f(e,e)\bigr).
\]
The prefactor is bounded by $C_N$ and
$d_o/\sqrt d\le C_N$.  Coercivity therefore gives
$|H_\Gamma|\le C_N(1+\sqrt{\mathcal T})$.
The Michael--Simon Inequality \cite{michaelsimon1973}, in the form
with an ordinary boundary integral stated explicitly by Brendle
\cite[Theorem~1]{brendle2021}, says
\[
 \left(\int_\Gamma\phi^{3/2}\,dV_\Gamma\right)^{2/3}
 \le C\left[
   \int_\Gamma(|\nabla_\Gamma\phi|+(1+|H_\Gamma|)\phi)\,dV_\Gamma
                     +\int_{\partial\Gamma}\phi\,dA\right]
\]
for nonnegative supported $\phi$. The graph over a compact smooth
patch is embedded by the fixed Nash embedding. Enlarging the ambient
Euclidean space if needed puts its codimension at least two, as in
Brendle's Theorem. That theorem applies first to $\phi+\varepsilon>0$;
letting $\varepsilon\downarrow0$ and using
$\sqrt{a^2+b^2}\le a+b$ gives the displayed form.
The boundary form follows
also by extending each smooth immersed graph slightly across its
boundary and letting a cutoff collar shrink: its derivative
integral tends to the boundary integral and its added curvature
integral tends to zero.  This is done for each smooth graph
before taking the uniform estimate.  On an inner face $f$ is
constant, so $dA=dA_g$.
Equation~\eqref{eq:fc-unweighted-trace} bounds that boundary
integral by
$C_N\int_\Gamma[(1+\sqrt{\mathcal T})\phi+
|\nabla_\Gamma\phi|]\,dV_\Gamma$.
Taking $\phi=|\omega|^4$ and using Cauchy--Schwarz proves
\begin{equation}\label{eq:fc-graph-sobolev}
 \|\omega\|_{L^6(dV_\Gamma)}^2
 \le C_N\int_\Gamma
       \bigl(|\nabla_\Gamma\omega|^2+
                             (1+\mathcal T)\omega^2\bigr)\,dV_\Gamma.
\end{equation}

For a compact base cutoff $\eta$ test the scalar equation by
$q=\eta^2e^{2kZ}/L$.  Its differential is
\[
 dq=L^{-1}e^{2kZ}
       [2\eta\,d\eta+\eta^2(2k\,dZ-(p+2)\,dt)].
\]
The principal exponential term contributes
$8k\eta^2e^{2kZ}|dZ|_{A_\chi}^2$ against the measure
$\sqrt D\,dV_g$.  Coercivity gives
\[
 4(p+2)|\langle dt,dZ\rangle_{A_\chi}|
       \le\tfrac14\delta_0\mathcal T+C_N|dZ|_{A_\chi}^2.
\]
For $k\ge k_*(N)$ the second term absorbs.  The drift terms obey
\begin{align*}
 2k|\langle K(w,\cdot),dZ\rangle_{A_\chi}|
      &\le k|dZ|_{A_\chi}^2+Ck|K|^2,\\
 (p+2)|\langle K(w,\cdot),dt\rangle_{A_\chi}|
      &\le\tfrac14\delta_0\mathcal T+C_N|K|^2.
\end{align*}
The cutoff terms have the same bounds with a remainder
$C_Nk e^{2kZ}(\eta^2+|d\eta|_g^2)$.  The penalty is bounded on
the entire exterior and contributes to this remainder.
Before the last inequality of
\eqref{eq:fc-boundary-flux-small}, the boundary term is at most
$C\eta^2e^{2kZ}\sqrt d\le C\eta^2e^{2kZ}$.
Using \eqref{eq:fc-unweighted-trace} with
$\phi=\eta^2e^{2kZ}$ bounds it by
\[
 C_N\int_\Gamma e^{2kZ}
  \bigl[\eta^2(1+\sqrt{\mathcal T}+k|\nabla_\Gamma Z|)
                                  +\eta|d\eta|_g\bigr]\,dV_\Gamma.
\]
Young's Inequality absorbs the $\sqrt{\mathcal T}$ term and the
term linear in the gradient; the latter costs $C_Nk$ times the
weighted volume.  After comparison of the graph measures we have
\begin{equation}\label{eq:fc-exponential-energy}
 \int_\Gamma e^{2kZ}\eta^2
               (\mathcal T+k|\nabla_\Gamma Z|^2)\,dV_\Gamma
 \le C_N k\int_\Gamma e^{2kZ}
                    (\eta^2+|d\eta|_g^2)\,dV_\Gamma,
 \quad k\ge k_*(N).
\end{equation}
Its constants are independent of $k$ and $R$.

For nested base patches $U_r\subset U_{r'}$, combine
\eqref{eq:fc-graph-sobolev} with
\eqref{eq:fc-exponential-energy} applied to $\eta e^{kZ}$:
\[
 \|e^Z\|_{L^{6k}(\Gamma|U_r)}
 \le\bigl[C_Nk^2(1+(r'-r)^{-2})\bigr]^{1/(2k)}
                \|e^Z\|_{L^{2k}(\Gamma|U_{r'})}.
\]
Iterate with $k_j=3^j k_*$ and geometrically decreasing gaps.
The sums $\sum k_j^{-1}$ and $\sum j/k_j$ converge, giving
\[
 S(r):=\sup_{U_r}e^Z
 \le C_N(r'-r)^{-A_N}
                 \|e^Z\|_{L^{2k_*}(\Gamma|U_{r'})}.
\]
Increase $k_*$ to at least two.  The local $L^2$ estimate
\eqref{eq:fc-graph-L2} and interpolation give
$S(r)\le C_N(r'-r)^{-A_N}S(r')^{1-1/k_*}$.
For any $\vartheta>0$, Young's Inequality turns this into
\[
 S(r)\le\vartheta S(r')+
                    C_{N,\vartheta}(r'-r)^{-A_Nk_*}.
\]
Use radii increasing geometrically to a fixed outer patch and
choose $\vartheta<2^{-A_Nk_*-1}$.  The resulting series converges.
The final remainder tends to zero because each individual smooth
solution has a finite supremum on that closed outer patch.  The
bound obtained is independent of this individual supremum.
A finite cover proves the lemma, including at inner faces.
\end{proof}

\begin{proposition}[Bounded variables in every stage]
\label{prop:fc-bounded-range}
There are $N_0$ and $R_{\min}(N)\to\infty$ such that, for
$N\ge N_0$ and $R\ge R_{\min}(N)$, every smooth solution of the
four-stage family with $t\ge-\epsilon$ satisfies
\[
 |h|\le C_0,\qquad |f|\le C_0/\tau,
 \qquad-\epsilon\le t\le Z\le C_N,
 \qquad |df|\le C_N.
\]
The constants are uniform in the homotopy and truncation.
\end{proposition}

\begin{proof}
It remains to extend the first-stage compact bound to the tail;
the tensor $K$ need not be compactly supported.  Everywhere, not
only at a critical point of $t$, differentiation gives
\begin{equation}\label{eq:fc-drift-identities}
 \nabla w=H^f A_d+p d\,dt\otimes w,
 \qquad
 d\log u=(p+2+pv)\,dt+H^f(w,\cdot).
\end{equation}
Also
\[
 A_\chi=I-\frac{4+p}{4+pv}w\otimes w,
 \qquad
 \partial_p\frac{4+p}{4+pv}=\frac{4d}{(4+pv)^2}.
\]
In differentiating $A_\chi$, the derivatives of $v$ have the form
$2a d\,H^f(\cdot,e)+2p d a^2dt$.  Thus every axial $H^f(e,e)$
or axial $dt$ term in
$u^{-1}\operatorname{div}_g(uA_\chi K(w,\cdot))$ has a factor
$d$ or $\chi$.  More explicitly its absolute value is bounded by
\[
 P_N|K|\bigl(
 |H^f_{TT}|+|H^f_{eT}|+d|H^f_{ee}|
          +|dt_T|+d|dt_e|\bigr)
       +P_Na|\nabla K|.
\]
For the $d\log u$ term this follows from
$A_\chi K(w,\cdot)(e)=\chi aK(e,e)$; for the other terms it
follows directly from \eqref{eq:fc-drift-identities} and the
displayed derivative of the scalar coefficient.  Coercivity and
Young's Inequality imply
\begin{equation}\label{eq:fc-drift-bound}
 \left|u^{-1}\operatorname{div}_g(uA_\chi K(w,\cdot))\right|
 \le\tfrac12\delta_0\mathcal T
                       +C_N(|K|^2+a|\nabla K|).
\end{equation}
Since $K=O(r^{-2})$, $\nabla K=O(r^{-3})$, and
$\tau a\sigma=\tau a^2\sqrt D/l\ge(\tau/e)a^2$, the second
term on the right is, sufficiently far out at fixed $N$, at most
\[
 \tfrac14\rho+\tfrac14\delta_0\tau a\sigma.
\]
Indeed the cross term is absorbed by a fixed fraction of
$\tau a^2$ with remainder $C_Nr^{-6}$, while
$r^{-4}=o(r^{-3-\delta})$ because $\delta<1$.
At an interior maximum of $Z$ above $M_0$, the principal
divergence contributes $4uA_\chi:\Hess Z\le0$.
Equations~\eqref{eq:fc-high-source} and
\eqref{eq:fc-drift-bound} contradict the scalar equation.
The compact upper bound and outer value zero therefore give a
global first-stage bound.

In stages two and three the drift is zero and the source is
positive above a fixed high level, so an interior maximum there
is impossible.  At an inner face
\[
 Z\le t+\tfrac18\log(1+e^2C_1^2/\tau^2).
\]
Consequently high boundary $Z$ implies $t>1$, where the boundary
multiplier makes the conormal derivative zero.  The boundary
point principle on the high superlevel excludes that maximum.
These maximum principles apply to each smooth solution with its
positive definite coefficients, so they require no prior uniform
ellipticity bound.  Stage four has $f=0$ and $t=Z$.  Its positive
scales have the same high-level argument; at scale zero its
homogeneous mixed problem gives $Z=0$.  This is uniform as the
scale tends to zero.

Finally $D=e^{8(Z-t)}$ and $l\ge\ell$ give
\[
 |df|^2=l^{-2}(e^{8(Z-t)}-1)
                     \le\ell^{-2}(e^{8(C_N+\epsilon)}-1).
\]
The height bound was already proved.  The finite set of compact
separation and collar choices, followed by the polynomial
absorption, fixes $N_0$.  Enlarge $R_{\min}(N)$ to contain all
fixed patches and the end thresholds.  The preceding arguments
apply to floor-contact solutions as well.
\end{proof}

\subsection{Classical compactness and the scalar solution operator}

The analytic estimates in Section~\ref{reg-section} are proved
independently of the deformation.  In particular they do not
assume a modulus of continuity for the axial coefficient.

\begin{proposition}[Classical bounds]\label{prop:fc-classical-bounds}
At fixed $\epsilon,N$, the solutions in
Proposition~\ref{prop:fc-bounded-range} have local bounds of every
finite classical order, up to the inner and outer faces, uniform
in $\xi$ and in all sufficiently large $R$.  On the end one can
use patches with a fixed positive radius and uniform constants.
\end{proposition}

\begin{proof}
On the bounded variable range $l,u,D$ are bounded above and away
from zero and $\chi\ge c_N>0$.  The normalized trace equation is
\[
 A_\chi:\Hess_g f=\frac{\sqrt D}{l}
                                      (F_\xi-\tr_{A_\chi}K).
\]
Its right side is bounded.  Theorem~\ref{reg-rank-one} gives
$Df$ H\"older and
\[
 \int_{B_r}|D^2f|^2\,dx\le C_Nr^{1+\beta_N}
 \quad\hbox{for some }\beta_N>0.
\]
The scalar principal coefficient in coordinates is
$4\sqrt{\det g}\,u A_\chi^{ij}$; the drift is bounded.
By the smooth implicit change $t=t(x,Z,Df)$,
$Dt=t_ZDZ+t_{Df}:D^2f$ plus controlled coordinate terms.
The source, including every homotopy modification, is therefore
bounded in absolute value by
$C_N(1+|DZ|^2+|D^2f|^2)$.
The prescribed inner flux is bounded and the outer value of $Z$
is zero.  Reflection contributes at most a bounded plane density,
which has mass $O(r^2)$.  Lemma~\ref{reg-natural-growth} now
gives a H\"older estimate for $Z$.

All smooth bounded-range coefficients of the trace equation are
then H\"older, so its Dirichlet Schauder estimate gives
$f\in C^{2,\theta}$.  At a face $Df$ is normal, and the scalar
condition solves for
\[
 \partial_\nu Z=G_\xi(x,Z,f,Df)
\]
with $G_\xi$ smooth on the bounded range: $A_\chi\nu=\chi\nu$.
The expanded scalar equation has H\"older leading matrix and
right side bounded by $C_N(1+|DZ|^2)$.
Proposition~\ref{reg-coupled-bootstrap} applies.  Its interpolation
after subtracting a constant uses small H\"older oscillation of
$Z$ to absorb the quadratic gradient term in local $W^{2,p}$
estimates.  The normal datum's trace norm is controlled by a
lower-order $W^{1,p}$ norm; the derivatives of $f$ entering that
datum are already bounded.  Subsequent Schauder estimates give
every finite order.  The data have uniform smooth end-patch
bounds, and the large outer spheres have uniform boundary-chart
constants.  The local supremum-of-patch-norm argument in
Proposition~\ref{reg-coupled-bootstrap} therefore makes all these
estimates independent of $R$.
\end{proof}

We first solve the trace equation for prescribed $Z$. This will define
the compact map used in the degree argument below, and the same scalar
result will also apply to the later spatially varying lapse.

\begin{proposition}[A scalar trace equation with bounded input]
\label{prop:fc-trace-solve}
Let $M$ be a compact smooth Riemannian three-manifold with smooth,
disjoint boundary components carrying constant Dirichlet values.
For prescribed $Z\in C^{1,b'}(M)$, $0<b'<1$, consider
\begin{equation}\label{eq:fc-general-trace}
 A_b(x,Z,Df):\Hess_g f=Q(x,Z,f,Df),
 \qquad A_b=I+(b-1)e\otimes e,
 \quad e=\nabla f/|df|.
\end{equation}
Suppose the actual coefficients extend smoothly at $df=0$, are
smooth on bounded ranges of their displayed variables, and the
principal coefficient is independent of the height $f$.
Suppose, uniformly on bounded input sets and any auxiliary compact
parameter family, the following hold:
\begin{enumerate}[label=(\roman*)]
 \item constant upper and lower height barriers bound $|f|\le M_1$
 for every solution and for the normalized interpolation to
 $\Delta_g f=c_*f$, with fixed $c_*>0$;
 \item on that height range, $0<b\le1$,
 $b\ge c/(1+|df|)$, and $|Q|\le C(1+|df|)$;
 \item $Q_f>0$, uniformly away from zero on bounded variable
 ranges.
\end{enumerate}
Then the Dirichlet equation has a unique solution in $C^{3,b'}$.
It depends continuously on the input and parameters and is bounded
in $C^{3,b'}$ on bounded input sets.  No prospective floor
condition on an implicit auxiliary variable is required.
\end{proposition}

\begin{proof}
Take a smooth extension of $M$ and a distance collar radius smaller
than a normal injectivity radius and half the separation of
distinct boundary components.  At a face use the barriers given
by its value plus or minus
\[
 \psi(r)=k^{-1}\log(1+Br),\qquad \psi''=-k(\psi')^2.
\]
At slopes $q=\psi'\ge1$, the upper barrier's axial contribution
is at most $-ckq/2$, whereas its tangential curvature and source
cost at most $C(1+q)$.  Choose $k$ large, then a collar radius
$r_0$ sufficiently small, and finally $B$ large with $Br_0\ge1$
and $\psi(r_0)>2\operatorname{osc}f$.
Then $q\ge1/(2kr_0)$ throughout the collar.  The lower barrier
has the reverse sign.  It suffices to check their residuals at
possible comparison contacts, where their heights are in the
fixed solution range; height monotonicity gives comparison.
These barriers exclude a positive maximum of
\[
 f(x)-f(y)-\psi(\operatorname{dist}(x,y))
 \quad\hbox{for }0<\operatorname{dist}(x,y)<r_0
\]
with one endpoint on $\partial M$.  Their value at $r_0$ excludes
the other edge of this two-point domain.

At a positive interior maximum, the endpoint gradients have
common length $q=\psi'(r)$ and are parallel transports along the
short joining geodesic.  Simultaneous endpoint variations in
parallel transverse directions, using the second variation of
distance, give
\[
 \tr_{e_x^\perp}\Hess f(x)-
 \tr_{e_y^\perp}\Hess f(y)\le Cqr.
\]
The transverse eigenvalues of both operators are exactly one.
Separate axial variations give
$f_{e_xe_x}(x)\le\psi''(r)$ and
$-f_{e_ye_y}(y)\le\psi''(r)$.  Therefore
\[
 -2C(1+q)\le Q(x)-Q(y)
 \le Cqr+(b_x+b_y)\psi''(r)
 \le Cqr-\frac{2ckq^2}{1+q},
\]
which contradicts the choices of $k$ and $r_0$.  No derivative or
continuity modulus of $b_x-b_y$ was used.  Interchanging the
points and then letting their distance tend to zero yields the
global Lipschitz bound $|df|\le\psi'(0)$.

Interpolate the normalized equation by
$A_s=(1-s)I+sA_b$ and $Q_s=(1-s)c_*f+sQ$.
The axial lower bound, source growth, constant barriers, and
strict height sign persist.  The preceding gradient estimate is
uniform in $s$, so the equations are uniformly elliptic.
Theorem~\ref{reg-rank-one} first gives $C^{1,\theta}$, then
Schauder gives $C^{2,\theta}$.  On this range $Df$ and the
prescribed $Z$ are Lipschitz; the undifferentiated coefficients
therefore have bounded $C^{b'}$ norms.  Schauder upgrades to
$C^{2,b'}$ and one further differentiation gives $C^{3,b'}$.
The linearization has positive definite principal matrix and
negative zeroth-order coefficient $-\partial_fQ_s$.
The Dirichlet maximum principle and linear elliptic theory give
its inverse.  The implicit function theorem supplies openness
of continuation from $\Delta_g f=c_*f$; the uniform estimates
and compactness supply closedness.  Comparison at a maximum of
the difference proves uniqueness, because the gradients agree
and the height source is strictly increasing.  The implicit
function theorem, or compactness and uniqueness, gives the
asserted dependence and uniform bounds.
\end{proof}

\begin{corollary}[The flat trace solution operator]
\label{cor:fc-flat-trace-solve}
For fixed $N\ge4$, $\tau>0$, a finite $\Omega_R$, and any
$\xi\in[0,4]$, the trace equation of
Proposition~\ref{prop:fl-homotopies} has a unique solution
$f=S_\xi(Z)\in C^{3,b'}$ for every $Z\in C^{1,b'}$ with zero outer
value.  It has the continuity and bounded-input estimates of
Proposition~\ref{prop:fc-trace-solve}, without a floor assumption.
\end{corollary}

\begin{proof}
The implicit definition \eqref{eq:fl-Z-change} is smooth and
strictly increasing in $t$ at fixed $df$.
The height estimates use only the trace equation, including its
strict height sign, so they hold without a floor.  For bounded
$Z$ and $\sigma\to\infty$, necessarily $t\to-\infty$,
$p=N$, and $l=e^{Nt}$.  The exact relation is
\[
 e^{8Z}=e^{8t}+\sigma^2e^{(2N+8)t}.
\]
Uniformly on bounded $Z$ ranges,
\begin{equation}\label{eq:fc-no-floor-asymptotics}
 t=\frac{4Z-\log\sigma}{N+4}+o(1),\qquad
 \chi\asymp\sigma^{-8/(N+4)},\qquad
 \frac{\sqrt D}{l}\asymp\sigma.
\end{equation}
Since $N\ge4$, the axial factor is at least $c/(1+\sigma)$.
After division by $l/\sqrt D$ the source is
\[
 Q=\frac{\sqrt D}{l}(F_\xi-\tr_{A_\chi}K),
 \qquad |Q|\le C(1+\sigma).
\]
All collar terms are bounded on the height range, and
$Q_f=(\sqrt D/l)\tau(F_\xi)_h>0$.
The principal matrix is independent of $f$ itself.  The normalized
interpolation to $\Delta f=\tau f$ retains these conditions and
the constant barriers.  Apply
Proposition~\ref{prop:fc-trace-solve}.
\end{proof}

\subsection{Degree, exhaustion, and end normalization}

\begin{lemma}[Degree on moving sections]\label{lem:fc-moving-degree}
Let $X$ be a Banach space and $H:[0,1]\times X\to X$ a continuous
map taking bounded sets to relatively compact sets.  Let
$\mathcal O\subset[0,1]\times X$ be a bounded relatively open
set.  If its closure has no fixed point on its relative product
boundary, then
$\deg(I-H_s,\mathcal O_s,0)$ is independent of $s$.
Here $\mathcal O_s=\{x:(s,x)\in\mathcal O\}$, and the degree of
an empty section is zero.
\end{lemma}

\begin{proof}
The fixed set in $\overline{\mathcal O}$ is compact by the joint
compactness and continuity, and it lies in $\mathcal O$ by
hypothesis.  For each parameter, surround its compact fixed-point
set by a bounded open neighborhood whose closure is contained in
all sufficiently nearby sections.  All nearby fixed points of
$\overline{\mathcal O}$ lie in this neighborhood: otherwise a
sequence outside it would converge to a fixed point at the
original parameter outside the neighborhood.  Excision identifies
the section degree with the degree on this one neighborhood, and
ordinary homotopy invariance makes it constant locally.  At a
parameter without fixed points the same compactness argument
makes all nearby degrees zero.  A finite covering of the
parameter interval proves the assertion.  These are the
compact-perturbation degree and excision properties of
Leray--Schauder \cite{lerayschauder1934}.
\end{proof}

\begin{proposition}[Finite-domain existence]\label{prop:fc-degree}
For each fixed sufficiently small $\epsilon>0$, all sufficiently
large $N$, and every $R\ge R_{\min}(N)$, the physical flat system
has a smooth solution on $\Omega_R$ with $t> -\epsilon$ on its
closure.  Its classical local bounds are uniform as $R\to\infty$
at fixed $N$.
\end{proposition}

\begin{proof}
Choose $N_0\ge\max\{4,\lceil2/\epsilon\rceil\}$ large enough
for Proposition~\ref{prop:fl-floor} and all the preceding
estimates.  Enlarge $R_{\min}(N)$ to tend to infinity and to
enforce the outer height and gradient barriers.  Explicitly the
outer slope has the bound
$\sigma\le C_{\mathrm{out}}\tau^{-1}R^{-1-\gamma}$.
It suffices also to require
\begin{equation}\label{eq:fc-outer-radius}
 R\ge\left(
 \frac{2C_{\mathrm{out}}\ell}
             {\tau\sqrt{e^{8\epsilon}-1}}\right)^{1/(1+\gamma)}.
\end{equation}
At outer $Z=0$ the implicit expression for $Z$ evaluated at
$t=-\epsilon$ is then strictly negative, so its monotonicity
excludes an outer floor contact.  The choices may depend
arbitrarily on fixed $N$; there is no restriction requiring a
polynomial outer radius.

Fix $N,R$ and let
\[
 X=\{Z\in C^{1,b'}(\overline{\Omega_R}):
                       Z|_{\partial_{\mathrm{out}}\Omega_R}=0\}.
\]
For an input $Z$, solve $f=S_\xi(Z)$ by
Corollary~\ref{cor:fc-flat-trace-solve} and compute all auxiliary
variables and source terms from this pair.  Denote the frozen
drift summand by $D_\xi$, the full frozen bulk source by
$E_\xi$, and the prescribed inner flux by $q_\xi$.
Define $H_\xi(Z)=\widetilde Z$ by
\begin{equation}\label{eq:fc-compact-map}
 \begin{cases}
 \operatorname{div}_g(4uA_\chi\nabla\widetilde Z+D_\xi)
                                      =E_\xi&\text{in }\Omega_R,\\
 (4uA_\chi\nabla\widetilde Z+D_\xi)_\nu=q_\xi&\text{on }B,\\
 \widetilde Z=0&\text{on the outer face}.
 \end{cases}
\end{equation}
Only the displayed principal gradient uses the output.
On bounded input sets, $4uA_\chi$ is positive definite and
$C^{1,b'}$, the drift is $C^{1,b'}$, the bulk source is $C^{0,b'}$,
and the boundary datum is $C^{1,b'}$.
In particular $Dt$ involves $DZ,D^2f$, and no second derivative
of the input $Z$.  The nonempty outer Dirichlet face gives the
Poincar\'e Inequality, so the mixed linear problem is uniquely
solvable by coercivity.  The inward-normal convention changes
only the sign of its boundary functional.  Linear regularity
gives $\widetilde Z\in C^{2,b'}$.  The compact embedding into
$X$ and continuous dependence make $H$ a continuous compact
homotopy on bounded sets, including its concatenation points.
Its fixed points satisfy the original equations and bootstrap
to smoothness.

After the trace solve set
\[
 \mathcal O=\{(\xi,Z):
        \min_{\overline{\Omega_R}}t_\xi[Z]>-\epsilon,
                                 \|Z\|_{C^{1,b'}}<M\}.
\]
This is relatively open because the trace solve and implicit
change are continuous.  Choose $M$ larger than the fixed-point
bound supplied by Propositions~\ref{prop:fc-bounded-range}
and~\ref{prop:fc-classical-bounds}, using their higher estimates
if necessary to reach the chosen exponent $b'$.
A fixed point in $\overline{\mathcal O}$ is smooth and has
$t\ge-\epsilon$, so these estimates exclude its norm boundary.
Equation~\eqref{eq:fc-outer-radius} and
Proposition~\ref{prop:fl-floor} exclude its floor boundary.
Thus Lemma~\ref{lem:fc-moving-degree}, with the parameter interval
rescaled, applies.  Notice that arbitrary inputs need not satisfy
the floor, and the map need not preserve $\mathcal O_\xi$.

At $\xi=4$ the unique trace solution is $f=0$, the drift vanishes,
and both bulk source and inner flux in
\eqref{eq:fc-compact-map} are zero for every input.  Its output
is identically zero.  The zero input has $t=0> -\epsilon$ and lies
inside the norm cutoff.  Hence the final degree is one, and so
is the physical degree.  A physical fixed point exists in
$\mathcal O_0$.  The asserted radius-independent local bounds
are Proposition~\ref{prop:fc-classical-bounds}.
\end{proof}

\begin{proposition}[Existence and end estimates]
\label{prop:fc-existence}
For every fixed sufficiently small $\epsilon>0$ and sufficiently large $N$, the
physical system has a smooth solution on $\Omega$, including its
inner boundary, with $t\ge-\epsilon$.  For some $c_N>0$ and
every fixed derivative order $j$,
\begin{equation}\label{eq:fc-end-decay}
 |\nabla^j f|\le C_{N,j}e^{-c_Nr},\qquad
 Z,t=O_2(r^{-1}).
\end{equation}
The differences $Z-t$ and $l^2df^2$, with their derivatives, decay
exponentially.  Moreover
\begin{equation}\label{eq:fc-flux-asymptotic}
 V=4\nabla_g t+O(r^{-3}),\qquad
 \mathcal F_N:=\lim_{r\to\infty}\int_{S_r}V_{\nu_g}\,dA_g
 \quad\hbox{exists and is finite}.
\end{equation}
\end{proposition}

\begin{proof}
Apply Proposition~\ref{prop:fc-degree} to any $R_i\to\infty$.
Uniform classical estimates and a diagonal subsequence give
smooth convergence on each compact subset, including the fixed
inner faces.  The equations and boundary conditions pass to the
limit, as does $t\ge-\epsilon$.  We prove the end bounds already
on the truncations so that their normalization does not escape
to infinity.

Where $C=0$ and $\tr_gK=0$, the trace equation is
\begin{equation}\label{eq:fc-tail-f-linear}
 A_\chi:\Hess_g f+B_f\cdot df-c_f f=0,
 \qquad c_f=\tau\frac{\sqrt D}{l},\qquad
 B_f=-\frac{4+p}{4+pv}\frac l{\sqrt D}K(\nabla f,\cdot).
\end{equation}
Indeed
$\tr_{A_\chi}K=-(1-\chi)K(e,e)$ and
\[
 \frac{\sqrt D}{l}(1-\chi)K(e,e)
        =\frac{4+p}{4+pv}\frac l{\sqrt D}K(\nabla f,\nabla f).
\]
At fixed $N$ the equation has uniformly elliptic principal
coefficient, bounded drift, and $c_f\ge\tau/e>0$.
For a sufficiently small $c_N>0$, the positive function
$\exp[-c_N(r-r_0)]$ is a supersolution of its negative-principal
operator beyond a fixed large $r_0$: the positive height term
dominates the $O(c_N^2)$ radial and $O(c_N/r)$ tangential Hessian
terms and the bounded drift term.  Multiply it to dominate the
inner sphere data; it also dominates the zero outer data.
Comparison for both signs gives exponential decay uniformly in
the truncation.  The already uniform smooth coefficient bounds,
followed by the homogeneous linear estimates for
\eqref{eq:fc-tail-f-linear}, give exponential decay of every
fixed derivative.  The corresponding zero-Dirichlet estimates
apply near the outer spheres.  Thus $Z-t$ and graph errors are
exponentially small with derivatives.

Using tracefreeness of $K$ on the tail, the quadratic form is now
\[
 \mathcal T=\tfrac12|K|^2+4(p(Z)+1)|dZ|^2+O(e^{-c_Nr}),
\]
and $A_\chi=I+O(e^{-c_Nr})$, $u=L(Z)+O(e^{-c_Nr})$.
The drift $uA_\chi K(w,\cdot)$ is exponentially small as well.
After decreasing $c_N$ to absorb polynomial factors, the scalar
equation becomes
\begin{equation}\label{eq:fc-tail-Z-equation}
\begin{aligned}
 \Delta_g Z+b_N(Z)|dZ|^2
 &=\tfrac14\rho-\tfrac14\Pi_N m_0(Z)\rho_0
                      +\tfrac18\delta_0|K|^2+O(e^{-c_Nr}),\\
 b_N(s)&=p(s)+2-\delta_0(p(s)+1).
\end{aligned}
\end{equation}
All fixed-order local derivatives of the exponential error are
controlled.  The right side is $O(r^{-3-\delta})$ on the bounded
solution range, even before the penalty disappears.
Define the increasing smooth function
\[
 \Phi_N(0)=0,\qquad
 \Phi_N'(s)=\exp\!\left(\int_0^s b_N(v)\,dv\right).
\]
Its derivative is bounded above and below on that range, and
$\Delta_g\Phi_N(Z)=O(r^{-3-\delta})$.
For $0<\delta'<\delta<1$,
\[
 -\Delta_g(r^{-1}-r^{-1-\delta'})
    =\delta'(1+\delta')r^{-3-\delta'}+O(r^{-4})
       \ge c r^{-3-\delta'}
\]
at large radius.  Comparison of both signs, with zero outer
values, gives $|\Phi_N(Z)|\le C_Nr^{-1}$ uniformly in truncation,
and hence $Z=O(r^{-1})$.

In the limit, sufficiently far out $t> -\epsilon+N^{-1}$, so the
penalty vanishes.  Rescale comparable-radius annuli to unit size.
The equation for $\Phi_N(Z)$ has a bounded rescaled source and
uniform smooth metric coefficients.  Local $W^{2,p}$ estimates
first give the scaled gradient bound.  Its source is a smooth
function of $Z$ times the symbol weights $\rho$, $|K|^2$, plus
the controlled exponential errors.  Schauder estimates then give
two symbol derivatives, proving \eqref{eq:fc-end-decay}.
Since $L(0)=1$, $u=1+O(r^{-1})$, and the graph and drift errors
are exponential, the first assertion of
\eqref{eq:fc-flux-asymptotic} follows.  The divergence $u\Xi$ is
integrable: its polynomial tail terms are $O(r^{-4})$ and
$O(r^{-3-\delta})$, and its remaining terms are exponential.
The divergence theorem on annuli proves the flux limit.
\end{proof}

\subsection{Curvature, full cut areas, and the mass loss}

\begin{lemma}[The comparison metric]\label{lem:fc-comparison-metric}
The solution of Proposition~\ref{prop:fc-existence} defines
\[
 \widehat g=e^{4t}(g+l^2df^2)
\]
as a smooth complete metric on $\Omega$, with $R_{\widehat g}\ge0$
and strictly negative mean curvature on $B$ in the normal into
$\Omega$.  Its end satisfies
\[
 \widehat g-\delta_{\mathrm{Eucl}}=O_2(r^{-1}),\qquad
 R_{\widehat g}=O(r^{-3-\delta}),\qquad
 E_{\widehat g}=E_g-\frac{\mathcal F_N}{8\pi}.
\]
\end{lemma}

\begin{proof}
At the physical endpoint $F=h+C$ and
$w(h)=\tau a\sigma$.  The scalar identity
\eqref{eq:fl-alg-identity} therefore gives
\begin{align*}
 \tfrac12e^{4t}R_{\widehat g}
 ={}&8\pi(\mu+J(w))+(1-\delta_0)\mathcal T
          +(1-\delta_0)\tau a\sigma\\
    &+w(C)-(h+C)\tr K-\rho
                     +m_0(t)\Pi_N(\rho_0+v\rho_1).
\end{align*}
The physical height range from
Lemma~\ref{lem:fc-collar-slopes}, $|w|\le1$, and the preparation
inequality for $C,\rho$ make this nonnegative.  That preparation
is needed only on the compact support of $\tr K$ and $dC$;
on the rest of the end it reduces to $\rho\le8\pi(\mu-|J|)$.
The boundary identity and prescribed flux give
$H+4\partial_\nu t=-N_B$.  Since $f$ is face-constant,
\[
 H_{\widehat g}=e^{-2t}\sqrt d\,(H+4\partial_\nu t)
                     =-e^{-2t}\sqrt d\,N_B<0.
\]
The pointwise inequality $\widehat g\ge e^{-4\epsilon}g$ proves
completeness, with the compact smooth boundary included.

Proposition~\ref{prop:fc-existence} gives the metric decay.
The conformal formula and exponential graph error give
\[
 R_{\widehat g}=e^{-4t}
       (R_g-8\Delta_g t-8|dt|^2)+O(e^{-c_Nr}).
\]
Equation~\eqref{eq:fc-tail-Z-equation} gives
$\Delta_g t=O(r^{-3-\delta})$, proving the required scalar
falloff.  Exponential graph errors have zero ADM contribution.
The conformal ADM formula and $t=O_2(r^{-1})$ give
\[
 E_{\widehat g}=E_g-\frac1{2\pi}
                     \lim_{r\to\infty}\int_{S_r}\partial_{\nu_g}t\,dA_g.
\]
Replacing the normal and area by their Euclidean versions costs
$o(1)$, and the $O(r^{-3})$ error in
\eqref{eq:fc-flux-asymptotic} has vanishing sphere integral.
This gives the stated normalization and sign.
\end{proof}

\begin{lemma}[Flux loss]\label{lem:fc-flux-loss}
For fixed $\epsilon>0$,
\[
 \mathcal F_N\ge-\varepsilon_N,
 \qquad 0\le\varepsilon_N\le
                 P_N(\ell+\ell^2/\tau)\longrightarrow0.
\]
The constants in this estimate do not use the fixed-$N$
classical regularity constants.
\end{lemma}

\begin{proof}
The inward normal on $B$ gives
\begin{equation}\label{eq:fc-integrated-flux}
 \mathcal F_N=\int_\Omega u\Xi\,dV_g+\int_B V_\nu\,dA_g.
\end{equation}
Both integrals are finite by the end estimates.
At a physical face, $V_\nu/u=-(1-a)H-N_Bd\ge-Cd$.
The signed slopes yield $\sqrt d\le C\tau/\ell$, so its negative
integral is at most $C(\tau/\ell)\int_B L$.
Apply \eqref{eq:fc-weighted-trace} with a constant test and a fixed
collar cutoff.  On those collars $\rho$ has a fixed positive
minimum, and therefore this loss is at most
\[
 P_N\frac\tau\ell
            \int_\Omega u(\rho+\delta_0\mathcal T)\,dV_g.
\]
For $N$ large it absorbs in one quarter of those positive terms.

On the penalty band $l,L\asymp\ell$, with constants independent
of $N$, and the exact definitions give
\[
 u\le C(\ell+\ell^2\sigma),\qquad
 uv\le C\ell^2\sigma,\qquad
 ua\sigma=e^{2t}l^2\sigma^2\ge c\ell^2\sigma^2.
\]
Thus, increasing only a polynomial constant,
\begin{align*}
 u m_0\Pi_N(\rho_0+v\rho_1)
 &\le P_N[\ell\rho_0+\ell^2\sigma(\rho_0+\rho_1)]\\
 &\le\tfrac14\delta_0\tau ua\sigma
       +P_N\left[\ell\rho_0+
                \frac{\ell^2}\tau(\rho_0^2+\rho_1^2)\right].
\end{align*}
The three remaining weight integrals are finite.  Their end
orders are $r^{-3-\delta}$, $r^{-6-2\delta}$, and
$r^{-4\gamma}$, respectively, with $4\gamma>3$.
After inserting both estimates in
\eqref{eq:fc-integrated-flux}, positive fractions of all three
bulk terms remain.  Dropping them proves the bound.
Finally $\ell^2/\tau=\tau/\ell=\ell^{1/2}$, which decays faster
than every fixed polynomial in $N$.
\end{proof}

\begin{theorem}[Completion of the flat deformation estimate]
\label{thm:fc-mass-completion}
The flat exterior of Theorem~\ref{thm:fl-deformation} satisfies
\[
 E_g\ge\sqrt{\frac{A_*}{16\pi}}.
\]
\end{theorem}

\begin{proof}
For each fixed large $N$, the metric of
Lemma~\ref{lem:fc-comparison-metric} has a strictly mean-negative
inner boundary and strictly mean-positive large coordinate
spheres.  Apply the three-dimensional trapped-region theorem
with zero tensor to obtain its full outermost minimal enclosing
frontier.  It is a compact smooth surface, possibly disconnected,
and bounds the connected exterior containing the end.
It is outer area-minimizing: its minimizing enclosure exists by
the direct method between the fixed barriers; minimal-obstacle
regularity and the maximum principle make that enclosure a
smooth minimal bounding frontier.  A strict enlargement would
contradict maximality of the full minimal trapped region.
These are also the outermost/outer-minimizing horizon conventions
in the Riemannian Penrose Theorem \cite{bray2001,braylee2009}.

Let this full cut be $\Gamma_N$.  Its exterior is complete with
boundary, has one AF end of order one through two derivatives,
and has nonnegative scalar curvature with pointwise falloff
$O(r^{-3-\delta})$.  These satisfy the boundary version of the
Riemannian Penrose Theorem \cite[Theorem~1.4]{braylee2009}; no
extension across the original boundary is needed.  Every cut in
the prepared exterior is a full enclosing cut of the original
flat obstacle.  The pointwise metric inequality and the
two-dimensional area scaling therefore give
\[
 |\Gamma_N|_{\widehat g}
        \ge e^{-4\epsilon}|\Gamma_N|_g
        \ge e^{-4\epsilon}A_*.
\]
Lemmas~\ref{lem:fc-comparison-metric} and
\ref{lem:fc-flux-loss} imply
\[
 E_g+\frac{\varepsilon_N}{8\pi}
       \ge E_{\widehat g}
       \ge\sqrt{\frac{|\Gamma_N|_{\widehat g}}{16\pi}}
       \ge e^{-2\epsilon}\sqrt{\frac{A_*}{16\pi}}.
\]
First the end was exhausted at fixed $N$.
Now let $N\to\infty$ with $\epsilon$ fixed, and then let
$\epsilon\downarrow0$.  Only numerical estimates are passed to
the last two limits; no smooth limit of the comparison metrics
as $N\to\infty$ is used.  This completes
Theorem~\ref{thm:fl-deformation}.
\end{proof}
